\chapter{Connectivité aux actions et aux options}\label{ch:nw:equities-and-options-connectivity}

Une firme qui veut couvrir l'ensemble du marché américain des options doit s'abonner à un flux consolidé réparti sur 96 lignes multicast, dimensionné par
son opérateur pour les rafales plutôt que pour les moyennes, et transmettre ses cotations par des ports dont elle ne peut dépasser
ni le nombre ni le budget de messages. L'opérateur du flux publie sa capacité en pics de 10 millisecondes~; la milliseconde la plus chargée de juillet 2026 a totalisé 0.35
million de messages, soit environ \qty{135}{\giga\bit\per\second} sur le fil pendant une milliseconde. Pour les options, la connectivité relève
d'abord de la planification de capacité, avant d'être une question de latence.

Le livre~1 a présenté les flux directs, les processeurs d'information sur les titres (SIP) et l'accès sponsorisé~; le livre~11, la cotation de masse, le
limiteur de messages et la Market Access Rule~; le livre~13, le feed handler. Ce chapitre les transforme en nombres de ports,
liens et cœurs~: les types de ports qu'une plateforme vend, la bande passante du flux consolidé d'options, la \omterm{def:nw:equities-and-options-connectivity:risk}{couche de risque
de l'accès sponsorisé}, qui s'intercale entre une firme et la plateforme, et l'ingénierie du trafic de cotations. Le plan est calculé par \texttt{firm.portplan}
à partir des documents datés de l'opérateur et des plateformes, avec une minute de forte activité simulée et signalée comme telle.

\section{Passerelles et types de ports}

\begin{definition}[Port de saisie d'ordres]\label{def:nw:equities-and-options-connectivity:port}
Un \emph{port de saisie d'ordres}\index{port de saisie d'ordres} est une connexion logique, dotée de sa propre session et de ses propres identifiants, par
laquelle un membre envoie des ordres, des cotations ou des annulations à la passerelle d'une plateforme et en reçoit les accusés de réception~; les plateformes vendent
des ports par type (saisie d'ordres, cotation groupée, purge ou annulation de masse, drop copy) et les facturent et les limitent port par port.
\end{definition}

La connectivité d'une plateforme comporte deux couches. La couche physique est celle du chapitre~9~: une \omterm{def:nw:colocation-products-and-how-they-are-sold:mmr}{interconnexion} (cross-connect), une vitesse de port, une
\omterm{def:nw:colocation-products-and-how-they-are-sold:cabinet}{baie de colocation}. La couche logique, ce sont les ports~: des sessions ouvertes sur la passerelle, chacune avec un type, des frais et des règles. Les plateformes d'options
séparent la cotation de la saisie d'ordres classique, car les cotations d'un teneur de marché représentent l'essentiel de leur trafic entrant.

La MIAX Express Interface illustre ce schéma. Un port Full Service accepte tous les messages, cotations groupées comprises~; un port Limited
Service accepte tous les messages sauf les cotations groupées~; un port Priority Mass Cancel ne traite que les annulations de masse, traitées avant les
autres flux de la firme. Chaque environnement d'appariement (un «~cloud~») négocie toutes les options sur un ensemble de sous-jacents, et une firme obtient
deux ports Full Service et douze ports Limited Service par cloud~: pour coter chaque sous-jacent, elle se connecte à chaque cloud. Sur
le Phlx de Nasdaq, l'interface de cotation est SQF, pour les cotations, les ordres exécuter-ou-annuler (immediate-or-cancel, IOC) et les réponses aux enchères~; le
dépôt réglementaire de la bourse indique qu'un seul port suffit à remplir les obligations de cotation d'un teneur de marché, et son barème de frais
plafonne un teneur de marché à 250 de ces ports.

\begin{dated}{2026-09}{Ports de cotation sur les bourses d'options américaines, d'après les dépôts et spécifications des plateformes}\label{dat:nw:equities-and-options-connectivity:ports}
\begin{itemize}
\item Phlx~: port SQF à 1\,185~\$ par port et par mois depuis janvier 2026, avec une réduction de 10~\% à 50~\% selon la part de volume du teneur de marché~; au plus 250
  ports SQF par teneur de marché. BOX~: 1\,000~\$ par mois pour l'ensemble des ports. MIAX~: frais MEI selon les classes cotées, jusqu'à 20\,500~\$ par mois
  pour plus de 100 classes. NYSE Arca et NYSE American~: 510~\$ par port pour les ports 1 à 40, 170~\$ à partir du port 41. Cboe EDGX Options~:
  750~\$ par port logique (tous ces montants tels que comparés dans le dépôt de Phlx de septembre 2025).
\item MIAX MEI~: 2 ports Full Service et 12 ports Limited Service par cloud~; un message Simple Bulk Quote transporte jusqu'à 50
  cotations unilatérales, avec 51 octets d'en-tête et 15 octets par cotation~; la spécification qualifie MEI de «~une interface
  de cotation synchrone~».
\end{itemize}
\end{dated}

Vingt-quatre ports coûtent 28\,440~\$ par mois sur le Phlx, 18\,000~\$ sur EDGX, 12\,240~\$ sur Arca ou American, 1\,000~\$ sur BOX, et
sur MIAX, le montant dépend des classes et non des ports (20\,500~\$ pour l'ensemble). Les drop copy et les ports de purge viennent en
sus. Les frais varient d'un facteur vingt-huit pour un même nombre de sessions~: ce nombre relève d'un choix de conception, et le barème
de frais est l'une de ses entrées.

\section{Bande passante du flux consolidé d'options}

\begin{definition}[Partition du flux]\label{def:nw:equities-and-options-connectivity:partition}
Une \emph{partition du flux}\index{partition du flux} est l'un des flux parallèles (lignes, canaux, \omterm{def:nw:networking-for-trading:group}{groupes multicast}) en lesquels
l'opérateur d'un flux répartit ses messages selon une règle publiée sur l'instrument, de sorte que chaque flux tienne sur une liaison et sur un handler, et qu'un abonné
ne retienne que les partitions dont il a besoin.
\end{definition}

Le flux consolidé d'options compte 96 lignes depuis son extension, et depuis le 1\up{er} juin 2026 une distribution publiée des symboles
répartit la charge entre elles~: par plage de symboles, et, au sein des symboles les plus actifs, par calls et puts et par mois ou jours d'échéance
impairs et pairs, «~pour un équilibrage optimal des symboles et une utilisation optimale de la capacité des lignes~». Son avis de capacité est rédigé pour
la planification des rafales (l'\cref{dat:nw:equities-and-options-connectivity:opra}).

\begin{dated}{2026-09}{Capacité du flux consolidé d'options, d'après les avis de son opérateur}\label{dat:nw:equities-and-options-connectivity:opra}
La projection en vigueur en juillet 2026 donne, pour un flux, 13.575 millions de messages, 4.403 gigabits et 1.564 million de paquets
pendant les 100 millisecondes les plus chargées~; 1.562 million de messages et 0.501 gigabit pendant les 10 plus chargées~; 311 milliards de messages par jour~;
et au plus 625\,000 messages par 100 millisecondes sur une ligne. Les deux flux redondants doublent la bande passante~; ajouter 10~\% pour
les retransmissions. La latence médiane à travers le processeur est inférieure à 18 microsecondes. Quatre-vingt-seize lignes, avec une nouvelle
distribution des symboles depuis le 1\up{er} juin 2026.
\end{dated}

Lu en débits, l'avis demande \qty{44.0}{\giga\bit\per\second} sur les 100 millisecondes les plus chargées et
\qty{50.1}{\giga\bit\per\second} sur les 10 plus chargées. Ses gigabits divisés par ses paquets donnent 352 octets par paquet et 8.68
messages par paquet~; en traitant ces octets comme charge utile UDP et en ajoutant l'Ethernet, l'IP, l'UDP et le surcoût propre à la ligne à l'aide de
\texttt{firm.netsim} donne 418 octets par paquet sur le fil, \qty{48.0}{bytes} par message. Sur le fil, le pic de 10 millisecondes
est de \qty{59.4}{\giga\bit\per\second} par flux, et, avec les deux flux et la marge pour les retransmissions,
soit \qty{130.6}{\giga\bit\per\second}.

Un port de commutateur ne moyenne pas sur 10 millisecondes. Pour planifier les liaisons, le chapitre simule une minute de forte activité en
intervalles de 1 milliseconde~: un niveau lentement variable (un log-AR(1) de persistance d'une seconde) multiplié par des pics rares (un intervalle sur mille) et un
léger bruit. Deux paramètres, la dispersion du niveau et l'amplitude des pics, sont ajustés sur deux rapports publiés~: la
milliseconde la plus chargée de juillet 2026 face aux 10 millisecondes les plus chargées, 3.93, et les 10 millisecondes les plus chargées de la projection face à
ses 100 millisecondes les plus chargées, 1.151. L'ajustement donne 3.90 et 1.152~; mis à l'échelle du pic de 10 millisecondes de la projection, la minute simulée
dont les 100 millisecondes les plus chargées contiennent 13.56 millions de messages, contre 13.575 dans l'avis, un contrôle que l'ajustement n'avait pas à réussir. Les
messages de chaque intervalle sont ensuite répartis entre les 96 lignes, avec un bruit lognormal autour de parts égales (dispersion 0.5, une hypothèse).

\omcode{code/firm/portplan/firm_portplan.py}{74}{82}{La minute simulée de forte activité~: un niveau lent multiplié par des pics rares, en intervalles de 1~ms.}\label{lst:nw:equities-and-options-connectivity:burst}

\begin{omfigure}
\begin{tikzpicture}
\begin{axis}[width=11.5cm,height=5.6cm,xmin=-300,xmax=300,ymode=log,ymin=8,ymax=400,
  xlabel={temps depuis la milliseconde la plus chargée (ms)},ylabel={un flux sur le fil (Gb/s, log)},
  tick label style={font=\footnotesize},label style={font=\small},grid=major,grid style={gray!25},
  legend style={font=\footnotesize,at={(0.5,-0.3)},anchor=north,legend columns=3},table/col sep=comma]
\addplot[omThm!70,thin] table[x=t_ms,y=gbps_1ms]{figdata/networks/22-equities-and-options-connectivity/burst.csv};
\addplot[omDef,thick] table[x=t_ms,y=gbps_10ms]{figdata/networks/22-equities-and-options-connectivity/burst.csv};
\addplot[omMeth,dashed] coordinates {(-300,10) (300,10)};
\legend{intervalles de 1~ms,moyenne sur 10~ms,une liaison de 10~Gb/s}
\end{axis}
\end{tikzpicture}

\omcaption{Simulation~: six cents millisecondes de la minute de forte activité autour de sa milliseconde la plus chargée, pour un flux, sur le fil. La
milliseconde la plus chargée atteint \qty{234}{\giga\bit\per\second}, et la moyenne sur 10 millisecondes qui l'entoure reste en dessous de
\qty{40}{\giga\bit\per\second}. Données~: \texttt{fig\_options.py}, \texttt{nw\_options.LINES\_GBPS}.}\label{fig:nw:equities-and-options-connectivity:burst}
\end{omfigure}

Les lignes sont ensuite rangées sur les liaisons, entières et dans l'ordre, et une liaison suffit à un percentile de planification lorsque
ce percentile de son débit sur 1 milliseconde tient dans \qty{10}{\giga\bit\per\second} moins la marge de 10~\%.

\omcode{code/firm/portplan/firm_portplan.py}{113}{121}{Liens pour un flux~: les groupes de lignes entières les moins nombreux dont le percentile de planification tient dans une liaison.}\label{lst:nw:equities-and-options-connectivity:links}

\begin{proposition}[Liaisons selon le percentile]\label{prop:nw:equities-and-options-connectivity:links}
Le nombre de liaisons dont un flux a besoin ne diminue jamais lorsque le percentile de planification augmente, et au 100\up{e}~percentile des intervalles de 1~ms
il vaut au moins la milliseconde la plus chargée du flux divisée par la capacité utile d'une liaison.
\end{proposition}

\begin{proof}
Un percentile plus élevé du débit de chaque groupe est au moins égal au percentile plus bas, donc tout rangement qui échoue au percentile plus bas échoue
au percentile plus élevé. Au maximum, chaque liaison transporte au plus sa capacité dans chaque intervalle, donc les liaisons transportent ensemble au plus leur
nombre multiplié par la capacité dans l'intervalle le plus chargé, lequel doit contenir la milliseconde la plus chargée de l'ensemble du flux.
\end{proof}

\begin{omfigure}
\begin{tikzpicture}
\begin{axis}[ybar,width=11cm,height=5.8cm,bar width=9pt,ymin=0,ymax=200,xtick=data,
  xticklabels from table={figdata/networks/22-equities-and-options-connectivity/links.csv}{pct},
  xlabel={percentile de planification des intervalles de 1~ms},ylabel={liaisons, les deux flux},
  nodes near coords,every node near coord/.append style={font=\scriptsize},
  tick label style={font=\footnotesize},label style={font=\small},grid=major,grid style={gray!25},
  legend pos=north west,legend style={font=\footnotesize},table/col sep=comma]
\addplot[fill=omThm!45,draw=omThm] table[x=x,y=links10]{figdata/networks/22-equities-and-options-connectivity/links.csv};
\addplot[fill=omDef!55,draw=omDef] table[x=x,y=links25]{figdata/networks/22-equities-and-options-connectivity/links.csv};
\legend{liaisons de 10 Gb/s,liaisons de 25 Gb/s}
\end{axis}
\end{tikzpicture}

\omcaption{Simulation~: liaisons nécessaires pour les deux flux du flux consolidé d'options, 96 lignes rangées entières, en fonction
du percentile des intervalles de 1~ms que le plan doit porter. Médiane~: 4~; 99\up{e}~: 12~; 99.99\up{e}~: 44~; la milliseconde la plus chargée~: 178 liaisons
de \qty{10}{\giga\bit\per\second} ou 32 de \qty{25}{\giga\bit\per\second}. Données~: \texttt{fig\_options.py},
\texttt{nw\_options.links\_table()}.}\label{fig:nw:equities-and-options-connectivity:links}
\end{omfigure}

La réponse dépend du percentile plus que de toute autre chose. À la médiane, quatre liaisons de \qty{10}{\giga\bit\per\second} transportent
les deux flux~; au 99\up{e}~percentile, douze~; au 99.99\up{e}, qui laisse six millisecondes de la minute au-dessus du budget,
quarante-quatre. La borne de la proposition pour la milliseconde la plus chargée donne 26 liaisons par flux (\qty{234}{\giga\bit\per\second}
divisés par \qty{9.09}{\giga\bit\per\second}), tandis que le rangement de lignes entières en demande 89, soit 178 pour les deux flux, presque une liaison par ligne, car
la milliseconde la plus chargée d'une seule ligne atteint \qty{6.8}{\giga\bit\per\second} et les rafales des lignes ne tombent pas dans la même milliseconde.
Avec des liaisons de \qty{25}{\giga\bit\per\second}, il en faut 32. Ce que le plan ne porte pas au niveau de la liaison, d'un
tampon de commutateur ou d'une perte est absorbé (chapitre~1), de sorte que le percentile est un choix sur la fréquence à laquelle la société accepte une file d'attente ou un trou.
La limite de lignes n'est pas la contrainte~: les 100 millisecondes les plus chargées de la ligne la plus active portent 171\,500 messages, soit un quart des
625\,000 que l'avis autorise.

Les gestionnaires constituent le second budget. À 10 millions de messages par seconde, un cœur (une hypothèse~: \qty{100}{\nano\second} par message, nettement au-dessus
du coût de décodage de quelques nanosecondes du livre~13, pour laisser de la place à la construction du carnet), le débit de pointe sur 10 millisecondes de 156.2 millions de messages par seconde
demande 16 cœurs par flux. Ils ne peuvent pas absorber la milliseconde la plus chargée, soit 609\,500 messages, sans file d'attente.

\omcode{code/networks/22-equities-and-options-connectivity/python/nw_options.py}{37}{44}{La file devant les cœurs des gestionnaires, sous forme de fluide sur les intervalles de 1~ms.}\label{lst:nw:equities-and-options-connectivity:backlog}

Avec 16 cœurs, la file atteint 449\,500 messages et met \qty{2.8}{\milli\second} à se vider~; avec 32, \qty{0.9}{\milli\second}~;
avec 48, \qty{0.27}{\milli\second}~; 64 cœurs ne font jamais la file. Les cœurs achetés pour le pic sur 10 millisecondes transforment la milliseconde la plus chargée
en millisecondes de carnets périmés, précisément au moment où le marché bouge.

\section{Couches de contrôle de risque sous accès sponsorisé}

\begin{definition}[Couche de risque de l'accès sponsorisé]\label{def:nw:equities-and-options-connectivity:risk}
Une \emph{couche de risque de l'accès sponsorisé}\index{couche de risque de l'accès sponsorisé} est le contrôle pré-négociation du courtier sponsor, implémenté en logiciel
ou dans une passerelle matérielle, par lequel passent tous les ordres et toutes les cotations d'une société sponsorisée en route vers la plateforme, et
qui ajoute sa latence de traitement à chaque message.
\end{definition}

Sous la Market Access Rule (livre~11), un courtier qui donne à une société l'accès à une plateforme sous son agrément doit contrôler
les messages de cette société avant qu'ils n'atteignent la plateforme. Le contrôle se situe sur le chemin, et sur une interface de cotation il se situe à chaque
aller-retour. Les fournisseurs publient la latence de leurs passerelles.

\begin{dated}{2026-09}{Latences publiées des passerelles de contrôle de risque matérielles}\label{dat:nw:equities-and-options-connectivity:risk}
La passerelle FPGA Risk Check de Magmio~: contrôles de risque et de conformité pré-négociation «~aussi bas que 200 nanosecondes~». Le FPGA
Pre-Trade Risk Check d'Algo-Logic~: contrôles «~nettement en dessous d'une microseconde~», pour la conformité à la Rule 15c3-5.
\end{dated}

Aucun courtier ne publie la latence de sa couche logicielle~; le chapitre suppose \qty{5}{\micro\second} dans chaque sens. Le coût d'une
couche de risque ne se limite pas aux microsecondes ajoutées à un message~: sur une interface synchrone, elle réduit ce qu'un port peut porter.

\section{Ingénierie du trafic de cotations}

La charge de cotations d'un teneur de marché est égale à ses classes multipliées par leurs séries, multipliées par la fréquence de mise à jour, multipliées par deux côtés. Les messages groupés portent
jusqu'à 50 cotations mono-face, de sorte que les messages sont cinquante fois moins nombreux que les cotations. Sur une interface synchrone, lue ici comme un
seul message groupé en vol par port jusqu'à son accusé de réception (une hypothèse à confirmer avec la plateforme), un port envoie un
message par aller-retour.

\omcode{code/firm/portplan/firm_portplan.py}{132}{139}{Ports~: un message par aller-retour sur chacun, et au moins un par moteur d'appariement.}\label{lst:nw:equities-and-options-connectivity:ports}

\begin{proposition}[Temps de rafraîchissement des cotations]\label{prop:nw:equities-and-options-connectivity:refresh}
Sur une interface de cotation synchrone à $k$ ports par moteur d'appariement, $b$ cotations par message et un aller-retour $\tau$ (couche de risque
incluse), le rafraîchissement de $Q$ cotations réparties uniformément sur $E$ moteurs prend $Q\tau/(Ekb)$~: linéaire en l'aller-retour et en
l'inverse du nombre de ports.
\end{proposition}

\begin{proof}
Chaque moteur reçoit $Q/E$ cotations en $Q/(Eb)$ messages~; ses $k$ ports envoient $k$ messages par aller-retour en parallèle, donc ils ont besoin
de $Q/(Ekb)$ allers-retours.
\end{proof}

\begin{omfigure}
\begin{tikzpicture}
\begin{axis}[width=11cm,height=5.6cm,xmin=0,xmax=10,ymin=0,ymax=42,
  xlabel={couche de risque, dans un sens (\unit{\micro\second})},ylabel={rafraîchissement complet (ms)},
  tick label style={font=\footnotesize},label style={font=\small},grid=major,grid style={gray!25},
  legend pos=north west,legend style={font=\footnotesize},table/col sep=comma]
\addplot[omThm,thick] table[x=risk_us,y=one_port]{figdata/networks/22-equities-and-options-connectivity/refresh.csv};
\addplot[omDef,thick] table[x=risk_us,y=two_ports]{figdata/networks/22-equities-and-options-connectivity/refresh.csv};
\legend{un port de cotation par moteur,deux ports de cotation par moteur}
\end{axis}
\end{tikzpicture}

\omcaption{Temps nécessaire pour rafraîchir les 600\,000 cotations de 1\,200 classes de 250 séries sur douze moteurs d'appariement, 50 cotations par
message, un message en vol par port, un aller-retour de \qty{20}{\micro\second} avant la couche de risque (hypothèses). Deux ports et
aucune couche de risque~: \qty{10}{\milli\second}~; une couche logicielle de \qty{5}{\micro\second}~: 15. Données~: \texttt{fig\_options.py},
\texttt{nw\_options.refresh\_curve()}.}\label{fig:nw:equities-and-options-connectivity:refresh}
\end{omfigure}

La firme du chapitre cote en moyenne 1\,200 classes de 250 séries, avec quatre mises à jour par seconde et par côté~: 2.4 millions de cotations par seconde,
48\,000 messages groupés, \qty{0.31}{\giga\bit\per\second} de saisie d'ordres. Avec un aller-retour de \qty{20}{\micro\second}, un port écoule
50\,000 messages par seconde~; le débit seul impose donc un port, les douze moteurs de la plateforme en imposent douze, et les deux ports de cotation
par moteur que la plateforme alloue en donnent vingt-quatre. La saisie d'ordres est contrainte par les messages et les moteurs, non par la bande passante~; le flux, lui, est
contraint par la bande passante.

Le cas contraignant est le mouvement de marché généralisé, celui où chaque cotation doit changer en même temps~: 600\,000 cotations. Deux ports par moteur les rafraîchissent
en \qty{10.0}{\milli\second}~; derrière la passerelle FPGA, \qty{10.2}{\milli\second}~; derrière la couche logicielle supposée,
\qty{15.0}{\milli\second}~; avec un port par moteur et sans couche, \qty{20.0}{\milli\second}. L'ingénierie du trafic de cotation consiste à
ne pas avoir besoin de ce rafraîchissement complet~: les cotations unilatérales permettent à la firme d'envoyer en premier les côtés importants (c'est le point
propre à la spécification), les ports de purge retirent les cotations plus vite qu'ils ne les mettent à jour, et l'ordre dans lequel les séries sont rafraîchies est une décision de risque.

\begin{method}[Dimensionner la connectivité d'un desk d'options]\label{met:nw:equities-and-options-connectivity:size}
\begin{enumerate}
\item Lire le dernier avis de capacité et les métriques du flux~; les convertir en bits sur le fil en tenant compte de l'en-tête de trame~; doubler pour les deux
  flux et ajouter la marge de retransmission.
\item Choisir un percentile de planification à l'échelle de la milliseconde, et répartir des partitions entières sur les liens en fonction de celui-ci~; décider ce que les
  tampons du commutateur absorbent au-delà.
\item Dimensionner les cœurs de traitement sur le pic de 10 millisecondes et vérifier la file de la milliseconde la plus chargée~; ajouter des cœurs jusqu'à ce que le délai soit
  acceptable.
\item À partir de la charge de cotation, calculer les messages, les ports par moteur et le temps de rafraîchissement complet, en incluant la couche de risque
  dans l'aller-retour.
\item Tarifer les ports selon la grille de chaque plateforme~; ajouter des drop copies et des ports de purge~; relancer le calcul à chaque nouvel avis de capacité.
\end{enumerate}
\end{method}

\begin{table}[ht]
\centering\small
\begin{tabular}{@{}llr@{}}
\toprule
Budget & Quantité & Plan \\
\midrule
Flux, un seul flux & pic de 10 ms sur le fil & \qty{59.4}{\giga\bit\per\second} \\
Flux, deux flux & avec 10~\% de retransmissions & \qty{130.6}{\giga\bit\per\second} \\
Liens & 99.99\up{e}~percentile, 10 Gb/s & 44 \\
Cœurs de traitement & pic de 10 ms, un seul flux & 16 \\
Milliseconde la plus chargée & file avec 16 cœurs & \qty{2.8}{\milli\second} \\
Cotation & cotations et messages par seconde & 2.4 millions~; 48\,000 \\
Ports de cotation & deux par moteur, douze moteurs & 24 \\
Rafraîchissement complet & deux ports, couche de risque logicielle & \qty{15.0}{\milli\second} \\
\bottomrule
\end{tabular}
\caption{Le plan de capacité du desk d'options du chapitre (flux issu de la projection de juillet 2026 et d'une simulation étiquetée~; cotation
à partir d'hypothèses énoncées). Données~: \texttt{nw\_options.feed\_plan()}, \texttt{quote\_plan()}.}\label{tab:nw:equities-and-options-connectivity:plan}
\end{table}

\begin{omfigure}
\begin{tikzpicture}[font=\small,
  box/.style={draw=omDef,fill=omDef!10,rounded corners,align=center,minimum height=9mm,inner sep=3pt},
  ven/.style={draw=omThm,fill=omThm!10,rounded corners,align=center,minimum height=9mm,inner sep=3pt}]
\node[ven] (opra) at (0,2.2) {flux consolidé\\96 lignes, 2 flux};
\node[box] (sw) at (3.9,2.2) {commutateur\\liens par percentile};
\node[box] (fh) at (7.6,2.2) {cœurs de traitement\\et file};
\node[box] (st) at (10.6,1.1) {stratégie};
\node[box] (rk) at (7.6,0) {couche de risque\\(courtier)};
\node[box] (pt) at (3.9,0) {ports de cotation\\par moteur};
\node[ven] (me) at (0,0) {moteurs\\d'appariement};
\draw[->,omThm] (opra) -- (sw); \draw[->,omThm] (sw) -- (fh); \draw[->,omThm] (fh) -| (st);
\draw[->,omDef] (st) |- (rk); \draw[->,omDef] (rk) -- (pt); \draw[->,omDef] (pt) -- (me);
\draw[->,omDef,dashed] (me.south) -- ++(0,-0.5) -| node[pos=0.25,below,font=\footnotesize] {accusés de réception~: un aller-retour par message groupé} (rk.south);
\end{tikzpicture}

\omcaption{Les deux budgets d'un desk d'options~: le côté flux est contraint par la bande passante et les cœurs de traitement, le côté cotation par
les ports, les moteurs et les allers-retours, avec la couche de risque du courtier sur chacun. Schéma.}\label{fig:nw:equities-and-options-connectivity:desk}
\end{omfigure}

\section{Tutoriel~: une table de capacité pour un desk d'options}\label{tut:nw:equities-and-options-connectivity:tut}

\begin{tutorial}
\textbf{Objectif.} À partir de l'avis de capacité du flux et d'une charge de cotation, calculer les liens, les cœurs et les ports.
\textbf{État final~:} le \cref{tab:nw:equities-and-options-connectivity:plan} et la \cref{fig:nw:equities-and-options-connectivity:links}.
\begin{tutsteps}
\item \textbf{Avis.} \texttt{firm\_portplan.projection} lit les lignes datées de juillet 2026~; \texttt{wire\_bytes\_per\_msg} ajoute
  l'en-tête de trame avec \texttt{firm.netsim}.
\item \textbf{Minute chargée.} \texttt{fit\_burst} ajuste le modèle de rafale aux deux ratios publiés~; \texttt{burst}
  (\cref{lst:nw:equities-and-options-connectivity:burst}) et \texttt{line\_split} donnent les débits de chaque ligne à la milliseconde.
\item \textbf{Liaisons et cœurs.} \texttt{links\_needed} (\cref{lst:nw:equities-and-options-connectivity:links}) pour chaque
  percentile~; \texttt{nw\_options.backlog} (\cref{lst:nw:equities-and-options-connectivity:backlog}) pour les cœurs.
\item \textbf{Ports.} \texttt{ports\_needed} (\cref{lst:nw:equities-and-options-connectivity:ports}) et
  \texttt{nw\_options.quote\_plan} avec chaque couche de risque.
\end{tutsteps}
\textbf{Pour aller plus loin.} Remplacer les parts égales par ligne par la distribution publiée des symboles et un modèle du trafic de chaque
symbole~; ajouter le tampon du commutateur avec \texttt{firm.netsim.egress} et compter les pertes à chaque percentile.
\end{tutorial}

\section{Construction~: le planificateur de ports et de bande passante}\label{bld:nw:equities-and-options-connectivity:portplan}

\begin{build}
\bfield{Rôle} Liens, cœurs et ports pour un desk qui reçoit un flux partitionné et cote par les ports de la plateforme~: les
  lignes actions et options du plan du chapitre~29.
\bfield{Interface} \texttt{firm\_portplan}~: \texttt{load\_capacity}, \texttt{projection}, \texttt{wire\_bytes\_per\_msg},
\texttt{burst}, \texttt{window\_peak}, \texttt{ratios}, \texttt{fit\_burst}, \texttt{line\_split}, \texttt{links\_needed},
\texttt{cores\_needed}, \texttt{quotes\_per\_s}, \texttt{port\_msgs\_per\_s}, \texttt{ports\_needed}, \texttt{bulk\_bytes}.
\bfield{Règles} Les capacités et les frais publiés sont des lignes datées avec leurs sources~; la minute chargée est une simulation étiquetée dont les
  paramètres libres sont ajustés sur des ratios publiés~; tout autre paramètre est une hypothèse énoncée.
\bfield{Tests d'acceptation} \texttt{code/firm/portplan/tests/}~: les lignes de capacité, un trafic plat donnant des ratios égaux à un, des lignes égales
  réparties à la main, et l'arithmétique des ports.
\bfield{Pour aller plus loin} Parts par ligne issues de la distribution publiée des symboles~; pertes par la \omterm{def:nw:networking-for-trading:egress}{file de sortie} de \texttt{firm.netsim}~; règles de
  port par plateforme.
\end{build}

\begin{omsources}
\item SIAC, projections de capacité OPRA (septembre 2025) et avis de distribution des symboles sur 96 lignes (mai 2026)~; métriques opérationnelles
  clés d'OPRA (août 2026).
\item SEC Release No.~34-103889 et Federal Register 90 FR 60788 (frais de port Phlx SQF)~; MIAX Options, MEI Interface
  Specification, version 2.10a.
\item Pages produits de Magmio et Algo-Logic pour les passerelles de risque FPGA.
\end{omsources}

\section{Exercices}

\begin{exercise}[$\star$]\label{exo:nw:equities-and-options-connectivity:1}
Quelles sont, à partir de la projection de juillet 2026, les 100 millisecondes et les 10 millisecondes les plus chargées, exprimées en débits en gigabits par seconde,
pour un flux, tel que l'avis le compte~?
\end{exercise}

\begin{exercise}[$\star$]\label{exo:nw:equities-and-options-connectivity:2}
Combien de messages en masse par seconde une firme envoie-t-elle pour 2.4 millions de cotations unilatérales par seconde sur MEI~?
\end{exercise}

\begin{exercise}[$\star$]\label{exo:nw:equities-and-options-connectivity:3}
Pourquoi une firme qui cote sur MIAX a-t-elle besoin de ports sur chacun des clouds, même si la charge est faible sur certains d'entre eux~?
\end{exercise}

\begin{exercise}[$\star\star$]\label{exo:nw:equities-and-options-connectivity:4}
En utilisant la \cref{prop:nw:equities-and-options-connectivity:refresh}, combien de temps prend un rafraîchissement complet avec deux ports par moteur derrière la
passerelle FPGA, et que donnerait un troisième port par moteur~?
\end{exercise}

\begin{exercise}[$\star\star$]\label{exo:nw:equities-and-options-connectivity:5}
Pourquoi le 99.99\up{e}~centile est-il une cible de planification exigeante à l'échelle de la milliseconde, et que se passe-t-il au-delà~?
\end{exercise}

\begin{exercise}[$\star\star$]\label{exo:nw:equities-and-options-connectivity:6}
Que coûtent vingt-quatre ports de cotation par mois sur chaque plateforme de l'\cref{dat:nw:equities-and-options-connectivity:ports}~?
\end{exercise}

\begin{exercise}[$\star\star\star$]\label{exo:nw:equities-and-options-connectivity:7}
\emph{Programmation.} Avec \texttt{nw\_options.backlog}, combien de cœurs de traitement maintiennent la file de la milliseconde la plus chargée sous
\qty{0.5}{\milli\second}~?
\end{exercise}

\begin{exercise}[$\star\star\star$]\label{exo:nw:equities-and-options-connectivity:8}
\emph{Trouver la faille.} «~L'avis indique \qty{50}{\giga\bit\per\second} au pic de 10 millisecondes, donc six
liens de \qty{10}{\giga\bit\per\second} par flux suffisent largement.~»
\end{exercise}

\section{Problème~: Cotation de toutes les séries}

\begin{problem}[{Problème du week-end --- ports et liens pour un teneur de marché d'options américain}]\label{pb:nw:equities-and-options-connectivity:1}
Une firme va coter en moyenne 1\,200 classes d'options de 250 séries chacune, avec quatre mises à jour par seconde et par côté, sur une plateforme dotée de douze
moteurs d'appariement et de cotation groupée synchrone, par l'intermédiaire d'un courtier parrain~; elle reçoit le flux consolidé dans son intégralité, les deux flux. Utiliser le
planificateur du chapitre.

\medskip\noindent\textbf{Partie I --- Le flux.}
\begin{enumerate}
\item Quelle bande passante la projection de juillet 2026 exige-t-elle, telle que l'avis la compte et telle qu'elle passe sur le fil, pour un seul flux et pour les deux avec la
  marge~?
\item Combien d'octets un message occupe-t-il sur le fil selon le chapitre, et à partir de quoi~?
\item Combien de liens de 10 Gb/s transportent les deux flux à la médiane, au 99\up{e}~centile et au 99.99\up{e}~centile des intervalles de 1 milliseconde~?
\item Pourquoi le transport de la milliseconde la plus chargée exige-t-il 178~?
\end{enumerate}

\medskip\noindent\textbf{Partie II --- Les handlers.}
\begin{enumerate}[resume]
\item Combien de cœurs la pointe de 10 millisecondes exige-t-elle à raison de 10 millions de messages par seconde et par cœur~?
\item Quelle file la milliseconde la plus chargée construit-elle avec eux, et combien de temps faut-il pour la vider~?
\item Combien de cœurs suppriment entièrement la file~?
\item Quel contrôle le modèle de rafale a-t-il réussi sans y avoir été ajusté~?
\end{enumerate}

\medskip\noindent\textbf{Partie III --- Les cotations.}
\begin{enumerate}[resume]
\item Combien de cotations, de messages en masse et de gigabits par seconde la firme envoie-t-elle~?
\item Combien de ports le débit exige-t-il, et combien la structure de la plateforme en impose-t-elle~?
\item Combien de temps prend un rafraîchissement complet avec deux ports par moteur, sans couche de risque, derrière la passerelle FPGA et derrière la
  couche logicielle~?
\item Que signifie le temps de rafraîchissement complet pour le risque de la firme après un mouvement de marché généralisé~?
\end{enumerate}

\medskip\noindent\textbf{Partie IV --- Le verdict.}
\begin{enumerate}[resume]
\item Énoncer le \emph{résultat nommé}~: les ports et les liens de 10 Gb/s nécessaires pour coter les classes de la firme au rythme de ses mises à jour et
  recevoir le flux consolidé dans son intégralité.
\item Quelles entrées sont publiées et lesquelles sont supposées~?
\item Que changeraient des liens de 25 Gb/s~?
\item Laquelle des décisions de la firme la couche de risque affecte-t-elle le plus~?
\item Dans quel ordre la firme doit-elle classer ses mises à jour de cotations après un mouvement~?
\item Que faut-il relancer lorsque le prochain avis de capacité arrive~?
\item Où ces chiffres prennent-ils place dans le plan du chapitre~29~?
\item En une phrase~: pourquoi la planification de la capacité de connectivité des options passe-t-elle avant la latence~?
\end{enumerate}
\end{problem}

\section{Questions d'entretien}

\begin{interviewq}[$\star$ \iqroles{développeur}]\label{iq:nw:equities-and-options-connectivity:1}
Pourquoi un opérateur de flux publie-t-il sa capacité par tranches de 10 millisecondes plutôt que par seconde~?
\end{interviewq}

\begin{interviewq}[$\star\star$ \iqroles{développeur}]\label{iq:nw:equities-and-options-connectivity:2}
Comment décideriez-vous du nombre de liens et de serveurs à acheter pour le flux complet des options américaines~?
\end{interviewq}

\begin{interviewq}[$\star\star$ \iqroles{développeur, chercheur}]\label{iq:nw:equities-and-options-connectivity:3}
Qu'est-ce qui limite la rapidité avec laquelle un teneur de marché peut mettre à jour l'ensemble de ses cotations, et comment la mesureriez-vous~?
\end{interviewq}

\begin{interviewq}[$\star\star$ \iqroles{développeur}]\label{iq:nw:equities-and-options-connectivity:4}
Un courtier parrain propose une couche de risque logicielle ou une passerelle matérielle. Que demanderiez-vous, et comment les compareriez-vous~?
\end{interviewq}

\begin{interviewq}[$\star\star$ \iqroles{trader}]\label{iq:nw:equities-and-options-connectivity:5}
Après un fort mouvement, quelles cotations faut-il mettre à jour en premier, et pourquoi~?
\end{interviewq}

\begin{interviewq}[$\star\star\star$ \iqroles{développeur, chercheur}]\label{iq:nw:equities-and-options-connectivity:6}
Concevez la connectivité d'un nouveau desk de market making d'options dans un centre de données américain~: flux, handlers, ports, couche de risque et
  surveillance.
\end{interviewq}
